\section{Introduction}
\label{sec:intro}

Language model agents act by calling tools, and a wrong call is rarely obvious when it is made: it parses,
its arguments have the right types, and the probabilities the model assigned to its tokens look ordinary, so
a transfer is made or a resource is provisioned before a person sees it. Whatever catches such an error must
act between the model producing the call and the system executing it, an interval too short to sample
continuations or consult a second model, and the only evidence available in it is the state of the model
that produced the call.

This paper asks how good a judge that state can be, and the question has a practical shape because
different signals need different access. Serving stacks expose token probabilities almost everywhere,
attention weights with some effort, and residual-stream activations rarely. A detector that needs
activations and one that needs only attention are not alternatives competing for one place in a system;
they belong to different deployments. We therefore report a frontier over levels of access rather than a
ranking of methods.

Two lines of work have proposed internal judges for this task: probes on the residual stream at the
structural positions of the call \citep{healy2026internal}, and spectral diagnostics of a layer of attention
read as a weighted graph \citep{noel2025gsp,binkowski2025lapeigvals}. Because they consume different tensors,
neither has served as a baseline for the other. We place them and the related detectors on the same data,
splits and labels, with every tool held out at test and every detector measured against baselines built from
lengths alone.

Three findings organise the results. First, whether an internal judge is worth building at all divides
cleanly by model family. On the families that hallucinate most, mean token log-probability is at chance or
below while probes on internal state are strong: these models make wrong calls confidently, which is the case
a guardrail exists for. On the more recent families we tested the picture reverses and confidence outscores
every trained detector; the obvious explanation, that a family which rarely fails starves its probe of
labels, fails a direct test. Second, where internals are needed, the residual-stream probe is the strongest
judge we measured and attention-only diagnostics come within a few points of it, so an operator without
activations is not reduced to reading confidence; both findings survive a conversation behind the call,
including one carrying a deliberately wrong earlier call. Third, how attention is read decides whether the
attention tier works at all: nearly every published diagnostic averages the attention matrix over heads
before computing anything, and analysing the heads separately moves the family from below a length baseline
to within reach of the residual stream, with width-matched controls attributing the gain to resolution.
Section~\ref{sec:theory} says why this must be so, and why the strongest published attention-only detector,
which already keeps heads apart and reads an in-degree sequence rather than a spectrum, is a baseline we
match and do not beat.

\paragraph{Contributions.}
\begin{enumerate}[topsep=2pt,itemsep=1pt,leftmargin=1.4em]
\item A controlled comparison of internal judges for hallucinated tool calls across five model families,
three attention designs and three data distributions, with a multi-turn condition, held-out tools, a length
floor in every table and a paired interval on every reported difference.
\item The finding that whether a trained judge beats output confidence is a property of the model family,
not of size, attention design or label budget, tested across runs and by thinning one run's positives.
\item An account of why head averaging limits attention-graph diagnostics: an annihilation result for any
averaged readout, a concavity result fixing the direction of the error, a measurement of how far real layers
sit from its equality condition, and controls separating resolution from width.
\item A negative result on combining the two families, which carry independent evidence yet do not combine
into a better detector at matched recall, and an audit of the measurement defects that make results in this
area fragile.
\end{enumerate}
